\documentclass[10pt,a4paper]{amsart}
\usepackage{amsmath,amssymb,amsthm,graphicx}
% For commutative diagrams we use Xy-pic
\usepackage[all]{xy}

% We use 2cell for 2-commutative diagrams.
\xyoption{2cell}
% Theorem environments.
%
\theoremstyle{plain}
\newtheorem{theorem}[subsection]{Theorem}
\newtheorem{proposition}[subsection]{Proposition}
\newtheorem{lemma}[subsection]{Lemma}

\theoremstyle{definition}
\newtheorem{definition}[subsection]{Definition}
\newtheorem{example}[subsection]{Example}
\newtheorem{exercise}[subsection]{Exercise}
\newtheorem{situation}[subsection]{Situation}

\theoremstyle{remark}
\newtheorem{remark}[subsection]{Remark}
\newtheorem{remarks}[subsection]{Remarks}

\numberwithin{equation}{subsection}

% Macros
%
\def\lim{\mathop{\mathrm{lim}}\nolimits}
\def\colim{\mathop{\mathrm{colim}}\nolimits}
\def\Spec{\mathop{\mathrm{Spec}}}
\def\Hom{\mathop{\mathrm{Hom}}\nolimits}
\def\Ext{\mathop{\mathrm{Ext}}\nolimits}
\def\SheafHom{\mathop{\mathcal{H}\!\mathit{om}}\nolimits}
\def\SheafExt{\mathop{\mathcal{E}\!\mathit{xt}}\nolimits}
\def\Sch{\mathit{Sch}}
\def\Mor{\mathop{\mathrm{Mor}}\nolimits}
\def\Ob{\mathop{\mathrm{Ob}}\nolimits}
\def\Sh{\mathop{\mathit{Sh}}\nolimits}
\def\NL{\mathop{N\!L}\nolimits}
\def\CH{\mathop{\mathrm{CH}}\nolimits}
\def\proetale{{pro\text{-}\acute{e}tale}}
\def\etale{{\acute{e}tale}}
\def\QCoh{\mathit{QCoh}}
\def\Ker{\mathop{\mathrm{Ker}}}
\def\Im{\mathop{\mathrm{Im}}}
\def\Coker{\mathop{\mathrm{Coker}}}
\def\Coim{\mathop{\mathrm{Coim}}}

% Boxtimes
%
\DeclareMathSymbol{\boxtimes}{\mathbin}{AMSa}{"02}

%
% Macros for moduli stacks/spaces
%
\def\QCohstack{\mathcal{QC}\!\mathit{oh}}
\def\Cohstack{\mathcal{C}\!\mathit{oh}}
\def\Spacesstack{\mathcal{S}\!\mathit{paces}}
\def\Quotfunctor{\mathrm{Quot}}
\def\Hilbfunctor{\mathrm{Hilb}}
\def\Curvesstack{\mathcal{C}\!\mathit{urves}}
\def\Polarizedstack{\mathcal{P}\!\mathit{olarized}}
\def\Complexesstack{\mathcal{C}\!\mathit{omplexes}}
% \Pic is the operator that assigns to X its picard group, usage \Pic(X)
% \Picardstack_{X/B} denotes the Picard stack of X over B
% \Picardfunctor_{X/B} denotes the Picard functor of X over B
\def\Pic{\mathop{\mathrm{Pic}}\nolimits}
\def\Picardstack{\mathcal{P}\!\mathit{ic}}
\def\Picardfunctor{\mathrm{Pic}}
\def\Deformationcategory{\mathcal{D}\!\mathit{ef}}
\usepackage{cleveref}
\crefname{originallemma}{original lemma}{original lemmas}
\Crefname{originallemma}{Original lemma}{Original lemmas}
\makeatletter
\newlabel{lemma-colimit}{{00D5}{source}{}{sourcepointer-1}{}}
\newlabel{lemma-colimit@cref}{{[originallemma][0][]00D5}{[1][0][]source}{}{}{}}
\makeatother
\crefname{originallemma}{original lemma}{original lemmas}
\Crefname{originallemma}{Original lemma}{Original lemmas}
\makeatletter
\newlabel{lemma-direct-sum-bounded-below-cohomology}{{0GIS}{source}{}{sourcepointer-2}{}}
\newlabel{lemma-direct-sum-bounded-below-cohomology@cref}{{[originallemma][0][]0GIS}{[1][0][]source}{}{}{}}
\makeatother
\crefname{originallemma}{original lemma}{original lemmas}
\Crefname{originallemma}{Original lemma}{Original lemmas}
\makeatletter
\newlabel{sites-cohomology-lemma-cohomology-modules-abelian-agree}{{03FD}{source}{}{sourcepointer-3}{}}
\newlabel{sites-cohomology-lemma-cohomology-modules-abelian-agree@cref}{{[originallemma][0][]03FD}{[1][0][]source}{}{}{}}
\makeatother
\crefname{originaltheorem}{original theorem}{original theorems}
\Crefname{originaltheorem}{Original theorem}{Original theorems}
\makeatletter
\newlabel{theorem-colimit}{{09YQ}{source}{}{sourcepointer-4}{}}
\newlabel{theorem-colimit@cref}{{[originaltheorem][0][]09YQ}{[1][0][]source}{}{}{}}
\makeatother
\crefname{originaltheorem}{original theorem}{original theorems}
\Crefname{originaltheorem}{Original theorem}{Original theorems}
\makeatletter
\newlabel{theorem-vanishing-affine-curves}{{03SC}{source}{}{sourcepointer-5}{}}
\newlabel{theorem-vanishing-affine-curves@cref}{{[originaltheorem][0][]03SC}{[1][0][]source}{}{}{}}
\makeatother
\title{Finiteness of etale cohomology of a constant sheaf with coefficients annihilated by an integer prime to the characteristic, or on a proper curve}
\author{The Stacks Project Authors}
\date{stacks-project:etale-cohomology.tex:local-snapshot-20261011; GNU FDL 1.2 or later}
\begin{document}
\maketitle
\noindent\textit{Editorial note.} All delivered mathematical text is the author's own, copied from the pinned Stacks Project chapter file etale-cohomology.tex; no statement or proof step was rewritten, reordered or removed.\par\medskip
\noindent\textit{Editorial note.} The counted proof is the author's complete argument: the proof is written out in full as a chain of reductions: the prime-factor induction, the change of the coefficient ring to Lambda/ellLambda, the passage to a finite free module by descending induction on q, and the basis decomposition of underline Lambda whose cohomology is computed by commutation with direct sums.\par\medskip
\noindent\textit{Editorial note.} The label lemma-constant-statements-coefficients names the counted statement, and the proof environment immediately following it in the source is delivered verbatim as its proof.\par\medskip
\noindent\textit{Editorial note.} The statement and proof are copied from the pinned chapter file; every reduction step and every cited input is retained exactly where the author placed it.\par\medskip
\noindent\textit{Editorial note.} The preamble is the chapter's own native preamble excerpt already used for this source: the Xy-pic setup (native preamble lines 14-18) and the native theorem declarations with the mathematical macros (native preamble lines 154-217).\par\medskip
\noindent\textit{Allowed standard prerequisites (editorial).}\begin{itemize}
\item Constant sheaf underline M on X\_etale associated to a Lambda-module. Same chapter, the constant sheaf on the etale site associated to a set, abelian group or module, and its description by the etale topology. Supplies the coefficient sheaves underline M, underline Lambda, underline F\_ell appearing throughout the reductions.
\item Long exact cohomology sequence and induction on the exponent. Same chapter, the long exact sequence of etale cohomology attached to a short exact sequence of etale sheaves. Supplies the exact sequences used to reduce from M annihilated by n to its prime factors and from M to a free module.
\item Vanishing above the dimension for the separated finite type curve. Same chapter, the vanishing and finiteness theorem for etale cohomology of tors/sheaves on separated finite type schemes of dimension <= 1 over an algebraically closed field. Supplies the vanishing in degrees > 2 used to run the descending induction on q, and the finiteness input of the same theorem.
\end{itemize}\medskip
\ifdefined\hypertarget\hypertarget{sourcepointer-1}{}\fi
\noindent\textit{Source reference.} Original lemma 00D5: sources/stacks-etale-cohomology/pinned-original-source.tex, source lines 7427--7427. Prerequisite pointer invoked at proof line(s) 15896; it names a standard fact of the same chapter that the counted argument uses, and every step of the counted proof built on it is retained verbatim in the delivered body.\par\medskip
\ifdefined\hypertarget\hypertarget{sourcepointer-2}{}\fi
\noindent\textit{Source reference.} Original lemma 0GIS: sources/stacks-etale-cohomology/pinned-original-source.tex, source lines 7831--7831. Prerequisite pointer invoked at proof line(s) 15897; it names a standard fact of the same chapter that the counted argument uses, and every step of the counted proof built on it is retained verbatim in the delivered body.\par\medskip
\ifdefined\hypertarget\hypertarget{sourcepointer-3}{}\fi
\noindent\textit{Source reference.} Original lemma 03FD: sources/stacks-chow/tags.txt, source lines 4488--4488. Prerequisite pointer invoked at proof line(s) 15854; it names a standard fact of a sibling chapter that the counted argument uses, and every step of the counted proof built on it is retained verbatim in the delivered body.\par\medskip
\ifdefined\hypertarget\hypertarget{sourcepointer-4}{}\fi
\noindent\textit{Source reference.} Original theorem 09YQ: sources/stacks-etale-cohomology/pinned-original-source.tex, source lines 7400--7400. Prerequisite pointer invoked at proof line(s) 15895; it names a standard fact of the same chapter that the counted argument uses, and every step of the counted proof built on it is retained verbatim in the delivered body.\par\medskip
\ifdefined\hypertarget\hypertarget{sourcepointer-5}{}\fi
\noindent\textit{Source reference.} Original theorem 03SC: sources/stacks-etale-cohomology/pinned-original-source.tex, source lines 15656--15656. Prerequisite pointer invoked at proof line(s) 15874; it names a standard fact of the same chapter that the counted argument uses, and every step of the counted proof built on it is retained verbatim in the delivered body.\par\medskip
\begin{lemma}
\label{lemma-constant-statements-coefficients}
Let $\Lambda$ be a Noetherian ring, let $M$ be a finite $\Lambda$-module which
is annihilated by an integer $n > 0$, let $k$ be an algebraically closed field,
and let $X$ be a separated, finite type scheme of dimension $\leq 1$ over $k$.
Then
\begin{enumerate}
\item $H^q_\etale(X, \underline{M})$ is a finite $\Lambda$-module
if $n$ is prime to $\text{char}(k)$,
\item $H^q_\etale(X, \underline{M})$ is a finite $\Lambda$-module
if $X$ is proper.
\end{enumerate}
\end{lemma}

\begin{proof}
If $n = \ell n'$ for some prime number $\ell$, then
we get a short exact sequence $0 \to M[\ell] \to M \to M' \to 0$
of finite $\Lambda$-modules and $M'$ is annihilated by $n'$.
This produces a corresponding short exact sequence of constant
sheaves, which in turn gives rise to an exact sequence
of cohomology modules
$$
H^q_\etale(X, \underline{M[n]}) \to
H^q_\etale(X, \underline{M}) \to
H^q_\etale(X, \underline{M'})
$$
Thus, if we can show the result in case $M$ is annihilated by
a prime number, then by induction on $n$ we win.

\medskip\noindent
Let $\ell$ be a prime number such that $\ell$ annihilates $M$.
Then we can replace $\Lambda$ by the $\mathbf{F}_\ell$-algebra
$\Lambda/\ell \Lambda$. Namely, the cohomology of $\mathcal{F}$
as a sheaf of $\Lambda$-modules is the same as the cohomology
of $\mathcal{F}$ as a sheaf of $\Lambda/\ell \Lambda$-modules,
for example by
Cohomology on Sites, Lemma
\ref{sites-cohomology-lemma-cohomology-modules-abelian-agree}.

\medskip\noindent
Assume $\ell$ be a prime number such that $\ell$ annihilates $M$
and $\Lambda$.
Let us reduce to the case where $M$ is a finite free $\Lambda$-module.
Namely, choose a short exact sequence
$$
0 \to N \to \Lambda^{\oplus m} \to M \to 0
$$
This determines an exact sequence
$$
H^q_\etale(X, \underline{\Lambda^{\oplus m}}) \to
H^q_\etale(X, \underline{M}) \to
H^{q + 1}_\etale(X, \underline{N})
$$
By descending induction on $q$ we get the result for $M$
if we know the result for $\Lambda^{\oplus m}$.
Here we use that we know that our cohomology groups
vanish in degrees $> 2$ by
Theorem \ref{theorem-vanishing-affine-curves}.

\medskip\noindent
Let $\ell$ be a prime number and assume that $\ell$ annihilates
$\Lambda$. It remains to show that the cohomology groups
$H^q_\etale(X, \underline{\Lambda})$ are finite $\Lambda$-modules.
We will use a trick to show this; the ``correct'' argument uses
a coefficient theorem which we will show later.
Choose a basis $\Lambda = \bigoplus_{i \in I} \mathbf{F}_\ell e_i$
such that $e_0 = 1$ for some $0 \in I$.
The choice of this basis determines an isomorphism
$$
\underline{\Lambda} = \bigoplus \underline{\mathbf{F}_\ell} e_i
$$
of sheaves on $X_\etale$. Thus we see that
$$
H^q_\etale(X, \underline{\Lambda}) =
H^q_\etale(X, \bigoplus \underline{\mathbf{F}_\ell} e_i) =
\bigoplus H^q_\etale(X, \underline{\mathbf{F}_\ell})e_i
$$
since taking cohomology over $X$ commutes with direct sums
by Theorem \ref{theorem-colimit}
(or Lemma \ref{lemma-colimit} or
Lemma \ref{lemma-direct-sum-bounded-below-cohomology}).
Since we already know that $H^q_\etale(X, \underline{\mathbf{F}_\ell})$
is a finite dimensional $\mathbf{F}_\ell$-vector space
(by Theorem \ref{theorem-vanishing-affine-curves}),
we see that $H^q_\etale(X, \underline{\Lambda})$
is free over $\Lambda$ of the same rank. Namely, given
a basis $\xi_1, \ldots, \xi_m$ of $H^q_\etale(X, \underline{\mathbf{F}_\ell})$
we see that $\xi_1 e_0, \ldots, \xi_m e_0$ form a $\Lambda$-basis for
$H^q_\etale(X, \underline{\Lambda})$.
\end{proof}

\end{document}
